\documentclass[11pt,a4paper]{article}
\usepackage[margin=25mm,headheight=14pt]{geometry}
\usepackage[T1]{fontenc}
\usepackage[utf8]{inputenc}
\usepackage{lmodern,amsmath,amssymb}
\usepackage{xcolor,graphicx,booktabs,longtable,array}
\usepackage{microtype,enumitem,fancyhdr,titlesec}
\usepackage{hyperref}
\definecolor{accent}{HTML}{163B55}
\hypersetup{colorlinks=true,linkcolor=accent,urlcolor=accent,
 pdftitle={Spectral type can vary within a limit-periodic Schrodinger hull},pdfauthor={Matthew J. Colbrook},
 pdfsubject={AIM 1018: mathematical result and proof}}
\urlstyle{same}
\Urlmuskip=0mu plus 1mu
\setlength{\parindent}{0pt}
\setlength{\parskip}{0.35em}
\setlength{\emergencystretch}{3em}
\setlist{itemsep=0.25em,topsep=0.3em}
\providecommand{\tightlist}{\setlength{\itemsep}{0pt}\setlength{\parskip}{0pt}}
\titleformat{\section}{\normalfont\large\bfseries\color{accent}}{\thesection}{0.7em}{}
\titleformat{\subsection}{\normalfont\normalsize\bfseries}{\thesubsection}{0.6em}{}
\pagestyle{fancy}
\fancyhf{}
\fancyhead[L]{\small\sffamily AIM 1018}
\fancyhead[R]{\small\sffamily Matthew J. Colbrook}
\fancyfoot[C]{\thepage}
\renewcommand{\headrulewidth}{0.3pt}
\setcounter{secnumdepth}{0}
\allowdisplaybreaks[2]
\title{\textbf{Spectral type can vary within a limit-periodic Schrödinger hull}}
\author{Matthew J. Colbrook\\[3pt]
{\small Department of Applied Mathematics and Theoretical Physics}\\
{\small University of Cambridge, United Kingdom}\\
{\small\href{mailto:mjc249@maths.cam.ac.uk}{\texttt{mjc249@maths.cam.ac.uk}}}}
\date{3 November 2026}
\begin{document}
\maketitle
\begin{abstract}
We construct a bounded real limit-periodic potential whose uniform hull contains both an operator with an explicit square-summable eigenfunction and an operator with no eigenvalues. The two phases have different nonzero spectral components, answering AIM 1018 negatively. The construction also permits arbitrarily small potentials.
\end{abstract}
\noindent\textbf{Repository target.}
This manuscript concerns AIM problem 1018 at revision
\texttt{3fa50884d129}. Its scope is stated below.
\medskip

\hypertarget{1018-status}{%
\section{Status}\label{1018-status}}

This is an elementary counterexample to the exact catalogue statement of
problem 1018, namely Damanik--Fillman Question~10.2
\cite[arXiv v1, p.~36]{DF20}; see also \cite[Remark~(v)]{DG11}.
Every construction and estimate needed below is included.
The construction gives an explicit phase with a positive,
square-summable eigenfunction at energy zero, and an explicit phase in
the same uniform hull with no eigenvalues at all. Consequently their
sets of nonzero spectral components differ, whether or not the first
operator has any additional spectral component.

\hypertarget{1018-periodic-triangular-functions-on-an-odometer}{%
\section{1. Periodic triangular functions on an
odometer}\label{1018-periodic-triangular-functions-on-an-odometer}}

Let \(G=\mathbb Z_2\) be the additive group of 2-adic integers, and
write \(x\mapsto x+1\) for its minimal translation. Set

\[
P_k=16^k,\qquad \epsilon_k=2^{-k},\qquad k\geq1,
\]

On the cyclic group \(\mathbb Z/P_k\mathbb Z\), use its usual
nearest-neighbor distance \(d_k\). Define the real, locally constant
function

\[
h_k(x)=\epsilon_k\,d_k\bigl(x\bmod P_k,0\bigr).
\tag{1}
\]

These functions satisfy:

\begin{enumerate}
\def\labelenumi{\arabic{enumi}.}
\tightlist
\item
  \(h_k\geq0\), \(h_k(-x)=h_k(x)\), and \(h_k(0)=0\).
\item
  \(h_k\) has period \(P_k\) under integer translation.
\item
  \(|h_k(x+1)-h_k(x)|\leq\epsilon_k\) for every \(x\in G\).
\end{enumerate}

Put

\[
b_k(x)=h_k(x+1)-h_k(x),\qquad
b(x)=\sum_{k=1}^\infty b_k(x).
\tag{2}
\]

Because \(\sum\epsilon_k=1\), this series converges uniformly to a
continuous real function with \(\|b\|_\infty\leq1\). Moreover,

\[
b(-x-1)=-b(x).
\tag{3}
\]

The functions \(h_k\) themselves need not have a convergent sum on all
of \(G\); no such claim is used.

\hypertarget{1018-a-bounded-even-sampling-function}{%
\section{2. A bounded, even sampling
function}\label{1018-a-bounded-even-sampling-function}}

Define

\[
F(x)=-e^{-b(x)}-e^{b(x-1)},\qquad
V_x(n)=F(x+n),\quad n\in\mathbb Z.
\tag{4}
\]

Then \(F\) is real and continuous, \(\|F\|_\infty\leq2e\), and (3) gives

\[
F(-x)=-e^{-b(-x)}-e^{b(-x-1)}
=-e^{b(x-1)}-e^{-b(x)}=F(x).
\tag{5}
\]

Let \(b^{(K)}=\sum_{k=1}^K b_k\) and
\(F^{(K)}(x)=-e^{-b^{(K)}(x)}-e^{b^{(K)}(x-1)}\). Each \(F^{(K)}\) is
even and \(P_K\)-periodic. The mean value theorem gives the uniform
bound

\[
\|F-F^{(K)}\|_\infty\leq2e\,2^{-K}.
\tag{6}
\]

Thus every \(V_x\) is a bounded uniform limit of periodic sequences,
exactly as required in the problem.

The map \(x\mapsto V_x\) is continuous from \(G\) to
\(\ell^\infty(\mathbb Z)\): this follows either from uniform continuity
of \(F\) in the translation-invariant group metric, or directly from (6)
and finite-level periodicity. Integer translations are dense in \(G\),
so

\[
\Omega_{V_0}=\{V_x:x\in G\}.
\tag{7}
\]

Injectivity of \(x\mapsto V_x\) is unnecessary.

\hypertarget{1018-an-explicit-eigenfunction-at-one-phase}{%
\section{3. An explicit eigenfunction at one
phase}\label{1018-an-explicit-eigenfunction-at-one-phase}}

For each integer \(n\), define

\[
H(n)=\sum_{k=1}^\infty h_k(n).
\tag{8}
\]

The series converges to a finite value: by the Lipschitz bound and \(h_k(0)=0\),

\[
0\leq h_k(n)\leq\epsilon_k|n|,
\qquad 0\leq H(n)\leq|n|.
\tag{9}
\]

Whenever \(P_k\geq2|n|\), the cyclic distance from \(n\) to zero equals
\(|n|\), so

\[
h_k(n)=\epsilon_k|n|.
\tag{10}
\]

For \(|n|\geq1\), choose the least \(K\geq1\) with \(16^K\geq2|n|\).
Then \(16^K\leq32|n|\), including the case \(K=1\). Therefore

\[
H(n)\geq |n|\sum_{k=K}^\infty2^{-k}
=2|n|\,2^{-K}
\geq \frac{2}{32^{1/4}}|n|^{3/4}.
\tag{11}
\]

Consequently

\[
u(n)=e^{-H(n)}
\tag{12}
\]

is positive and belongs to \(\ell^2(\mathbb Z)\). It is nonzero because
\(u(0)=1\). Both \(H\) and \(u\) are even sequences. The absolutely
convergent series in (8) and uniformly convergent series in (2) show

\[
H(n+1)-H(n)=b(n).
\tag{13}
\]

Thus

\[
\frac{u(n+1)}{u(n)}=e^{-b(n)},\qquad
\frac{u(n-1)}{u(n)}=e^{b(n-1)}.
\]

Substituting (4) gives the exact eigenvalue equation

\[
u(n-1)+u(n+1)+V_0(n)u(n)=0.
\tag{14}
\]

Hence \(0\) is an eigenvalue of \(H_{V_0}\), and
\(\mathrm{pp}\in\mathcal T(V_0)\).

For context, this eigenvalue is at the upper spectral edge. For every
finitely supported sequence \(f\) and every phase \(x\),

\[
-\langle f,H_{V_x}f\rangle
=\sum_{n\in\mathbb Z}e^{b(x+n)}
\left|f(n+1)-e^{-b(x+n)}f(n)\right|^2\geq0.
\]

This follows by expansion; boundedness extends it to all of \(\ell^2\).
Thus every operator is nonpositive, while the phase zero operator has
the displayed nonzero kernel vector. The potential \(V_0\) is
necessarily nonperiodic: a full-line periodic second-order difference
operator has no square-summable eigenfunction, as follows directly by
iterating its fixed, invertible period transfer matrix in both
directions.

\hypertarget{1018-an-explicit-phase-with-arbitrarily-accurate-distant-reflections}{%
\section{4. An explicit phase with arbitrarily accurate distant
reflections}\label{1018-an-explicit-phase-with-arbitrarily-accurate-distant-reflections}}

Set \(K_1=1\). Recursively define

\[
n_j=\sum_{r=1}^j P_{K_r},\qquad
K_{j+1}=n_j^2+K_j+1.
\tag{15}
\]

The sequence \(K_j\) is strictly increasing, and

\[
\theta_*=-\sum_{r=1}^\infty P_{K_r}
\tag{16}
\]

converges in \(G\). For each \(j\), the element \(\theta_*+n_j\) is
divisible by \(P_{K_{j+1}}\). By evenness and periodicity of
\(F^{(K_{j+1})}\),

\[
F^{(K_{j+1})}(\theta_*+2n_j-m)
=F^{(K_{j+1})}(\theta_*+m),\qquad m\in\mathbb Z.
\]

Writing \(W=V_{\theta_*}\) and using (6), we get

\[
d_j:=\sup_{m\in\mathbb Z}|W(2n_j-m)-W(m)|
\leq4e\,2^{-K_{j+1}}
\leq4e\,2^{-n_j^2}.
\tag{17}
\]

In particular \(n_j\to+\infty\) and \(e^{C n_j}\sqrt{d_j}\to0\) for
every fixed \(C>0\).

\hypertarget{1018-elementary-reflection-lemma-excluding-all-eigenvalues}{%
\section{5. Elementary reflection lemma excluding all
eigenvalues}\label{1018-elementary-reflection-lemma-excluding-all-eigenvalues}}

\textbf{Lemma.} Let \(W\) be a bounded real sequence. Suppose there are
positive integers \(n_j\to\infty\) and nonnegative \(d_j\) satisfying

\[
\sup_m|W(2n_j-m)-W(m)|\leq d_j,
\qquad e^{C n_j}\sqrt{d_j}\longrightarrow0
\]

for every \(C>0\). Then \(H_W\) has no \(\ell^2(\mathbb Z)\)
eigenfunction.

\textbf{Proof.} Suppose \(H_Wu=Eu\) with \(u\) real, nonzero, and
\(\|u\|_2=1\). Put \(v_j(m)=u(2n_j-m)\) and

\[
\mathcal W_j(m)=u(m+1)v_j(m)-u(m)v_j(m+1).
\]

Subtracting the two difference equations yields

\[
\mathcal W_j(m)-\mathcal W_j(m-1)
=[W(2n_j-m)-W(m)]u(m)v_j(m).
\]

Since both sequences tend to zero at infinity, summation from
\(-\infty\) and Cauchy--Schwarz give \(|\mathcal W_j(m)|\leq d_j\). At
the reflection center,

\[
|u(n_j)[u(n_j+1)-u(n_j-1)]|\leq d_j.
\tag{18}
\]

If \(|u(n_j)|\leq\sqrt{d_j}\), set \(w_j=u+v_j\). Then

\[
w_j(n_j)=2u(n_j),\qquad
w_j(n_j+1)=(E-W(n_j))u(n_j),
\]

so this two-component vector is \(O_{E,\|W\|_\infty}(\sqrt{d_j})\).
Otherwise set \(w_j=u-v_j\). Then \(w_j(n_j)=0\) and (18) gives
\(|w_j(n_j+1)|\leq\sqrt{d_j}\).

In either case, \((H_W-E)w_j\) has absolute value at most \(d_j\) at
each site. The inverse one-step transfer matrices have norm at most a
constant \(C_E\geq1\) depending only on \(E\) and \(\|W\|_\infty\).
Propagating the two-component vector backwards from \(n_j\) to zero
gives

\[
|w_j(0)|+|w_j(1)|
\leq C'_E C_E^{n_j}(\sqrt{d_j}+n_jd_j)\longrightarrow0.
\]

But \(v_j(0)=u(2n_j)\to0\) and \(v_j(1)=u(2n_j-1)\to0\), so necessarily
\(u(0)=u(1)=0\). The second-order recurrence then forces \(u\equiv0\), a
contradiction. This also covers \(d_j=0\) by the same case distinction.
\(\square\)

Applying the lemma to (17) shows
\(\mathrm{pp}\notin\mathcal T(V_{\theta_*})\).

\hypertarget{1018-conclusion-small-potentials-and-limits-of-the-assertion}{%
\section{6. Conclusion, small potentials, and limits of the
assertion}\label{1018-conclusion-small-potentials-and-limits-of-the-assertion}}

Both \(V_0\) and \(V_{\theta_*}\) lie in the same uniform limit-periodic
hull by (7), but

\[
\mathrm{pp}\in\mathcal T(V_0),\qquad
\mathrm{pp}\notin\mathcal T(V_{\theta_*}).
\]

Therefore the answer to the exact problem 1018 is \textbf{no}.

\hypertarget{1018-both-behaviors-occur-densely-in-the-hull}{%
\subsection{Both behaviors occur densely in the
hull}\label{1018-both-behaviors-occur-densely-in-the-hull}}

The hull elements with a nonzero point component include the translation
orbit of \(V_0\), which is dense. The elements with no point spectrum
contain a dense \(G_\delta\) subset, as follows directly from the
reflection lemma.

For an integer \(n\), write \((R_nW)(m)=W(2n-m)\). In the compact metric
hull \(\Omega=\Omega_{V_0}\), define

\[
\mathcal U_j=\bigcup_{n\geq j}
\{W\in\Omega:\|W-R_nW\|_\infty<e^{-n^2}\},\qquad j\geq1.
\]

These sets are open because each reflection map is an isometry in the
uniform norm. Each \(V_{-n}\) is exactly symmetric about \(n\), since
\(F\) is even. The set \(\{V_{-n}:n\geq j\}\) is dense in the hull, so
\(\mathcal U_j\) is dense. Baire's theorem makes
\(\bigcap_{j\geq1}\mathcal U_j\) a dense \(G_\delta\). Every potential
in that intersection has an unbounded sequence of reflection centers
with errors at most \(e^{-n^2}\), and hence has no eigenvalues by the
lemma. Thus the two different point-spectrum behaviors both occur
densely.

\hypertarget{1018-arbitrarily-small-examples}{%
\subsection{Arbitrarily small
examples}\label{1018-arbitrarily-small-examples}}

The counterexamples can be made arbitrarily small after a constant
energy shift. Indeed, replace \(\epsilon_k\) by \(\delta2^{-k}\) with
any \(\delta>0\). Then \(\|b\|_\infty\leq\delta\), the lower bound (11)
is multiplied by \(\delta\), and the eigenfunction remains
square-summable. Moreover,

\[
\|F+2\|_\infty\leq2(e^\delta-1)\longrightarrow0
\quad(\delta\downarrow0).
\]

Adding the constant \(2\) to every hull element simply shifts its
spectrum by \(2\), preserving all spectral types. The reflection error
bound becomes \(4e^\delta\delta2^{-K_{j+1}}\), which still satisfies the
lemma. Thus hull spectral-type variation occurs for arbitrarily small
uniform limit-periodic potentials.

The proof does not classify the rest of either operator's spectrum, does
not show that \(H_{V_0}\) is pure point, and by itself does not settle
the separate mixed-type question 1021 or the coupling question 1020. It
is unnecessary to classify those remaining components to distinguish the
two sets \(\mathcal T\).

The construction is reproducible from the displayed formulas alone and
uses no numerical evidence. The related classical theory of continuous
cocycles with discrete orbits over odometers provides conceptual
context, but is not invoked in the proof: J. Kwiatkowski and A.
Siemaszko, \emph{Discrete orbits in topologically transitive cylindrical
transformations}, Discrete Contin. Dyn. Syst. \textbf{27} (2010),
945--961, \url{https://doi.org/10.3934/dcds.2010.27.945}.

\begin{thebibliography}{3}
\raggedright
\bibitem{DF20}
D. Damanik and J. Fillman,
\emph{Spectral Properties of Limit-Periodic Operators},
in \emph{Analysis and Geometry on Graphs and Manifolds},
M. Keller, D. Lenz and R. K. Wojciechowski (eds.),
London Mathematical Society Lecture Note Series~461,
Cambridge University Press (2020), 382--444.
\href{https://doi.org/10.1017/9781108615259.016}{doi:10.1017/9781108615259.016}.
The question and section locators used here refer to
\url{https://arxiv.org/abs/1802.05794v1}.
\bibitem{DG11}
D. Damanik and Z. Gan,
\emph{Spectral properties of limit-periodic Schr\"odinger operators},
Communications on Pure and Applied Analysis \textbf{10} (2011),
859--871.
\href{https://doi.org/10.3934/cpaa.2011.10.859}{doi:10.3934/cpaa.2011.10.859}.
\url{https://arxiv.org/abs/0906.3337v3}.
\bibitem{KS10}
J. Kwiatkowski and A. Siemaszko,
\emph{Discrete orbits in topologically transitive cylindrical transformations},
Discrete and Continuous Dynamical Systems \textbf{27} (2010), 945--961.
\href{https://doi.org/10.3934/dcds.2010.27.945}{doi:10.3934/dcds.2010.27.945}.
This reference supplies context for discrete cocycles, not a theorem used
in the self-contained construction.
\end{thebibliography}
\end{document}
